\documentclass[11pt]{article}
\usepackage[margin=0.9in]{geometry}
\usepackage{amsmath,amssymb,amsthm,booktabs}
\usepackage[T1]{fontenc}
\usepackage{lmodern,microtype}
\usepackage[colorlinks=true,linkcolor=blue,urlcolor=blue]{hyperref}
\newtheorem{theorem}{Theorem}[section]
\newtheorem{lemma}[theorem]{Lemma}
\newtheorem{proposition}[theorem]{Proposition}
\newtheorem{corollary}[theorem]{Corollary}
\theoremstyle{remark}
\newtheorem{remark}[theorem]{Remark}
\newcommand{\norm}[1]{\left\|#1\right\|}
\title{Toward a four-dimensional Yang--Mills construction\\
\large Notebook 14: nonflat backgrounds and energy stability}
\author{Research working document}
\date{16 October 2026}
\begin{document}
\maketitle
\begin{abstract}
We pass from fixed-reference nonlinear averages to the actual
moving-reference averaging operators in a small global connection
chart. An explicit endpoint term corrects the change of derivative
coordinates. An exact cumulative gauge relation controls ordinary
reference backgrounds uniformly through the blocking depth, giving
small perturbations of both source averaging operators in physical
weighted norms. A separate Wilson plaquette calculation and a form
resolvent argument then control a gauge-fixed inverse, including the
scalar compression that defines its gauge projection. The inverse
estimates concern the energy blocks, with at most one derivative on
each side. They apply on a small global physical C1 ball, not yet on
the source's full local-frame and retained-history domains.
\end{abstract}
\noindent\textbf{Status.}
Research proofs checked through separate internal derivations and
audits. No external certification or novelty claim is made. The
explicit source-coordinate matches below refer to the deposited
transcriptions at the pinned commit, not to newly recovered original
CMP98 or CMP99b papers. The action and transpose normalizations are
stated explicitly at the point where the Wilson operator is used.